% GENERATED FILE. Edit canonical lessons, then run npm run generate:books.
% canonical-source-hash: fnv1a64:0adb169259648816
% canonical-lessons: ML-C420-pracaranam, ML-C420-prakatanapatrika, ML-C420-bhuripaksam, ML-C420-nyunapaksam, ML-C420-vijayikkuka, ML-R420-review-413-420

\chapter{Winning a Seat}
\label{ch:ml-winning-a-seat}
\hlchaptermodality{420}

\begin{chapteropening}
\textbf{By the end of this chapter:} \emph{I can say a campaign, a manifesto, a majority, a minority and to win.}

Complete the last lesson of chapter 420: I can say a campaign, a manifesto, a majority, a minority and to win.
\end{chapteropening}

\section[\texorpdfstring{\ml{പ്രചാരണം} (pracāraṇaṁ)}{പ്രചാരണം (pracāraṇaṁ)}]{\ml{പ്രചാരണം} (pracāraṇaṁ) — a campaign}

\label{lesson:ML-C420-pracaranam}



\begin{quote}
Before the new one: say the Malayalam for a candidate, then the Malayalam for a ballot.
\end{quote}

\subsection*{You'll want to know: \ml{പ്രചാരണം}}
\textbf{\ml{പ്രചാരണം}} — \emph{pracāraṇaṁ} — \textquotedblleft{}a campaign\textquotedblright{}.

\begin{quote}
\ml{തിരഞ്ഞെടുപ്പ്} \ml{പ്രചാരണം} \ml{ഇന്നലെ} \ml{അവസാനിച്ചു}.
\end{quote}

\begin{quote}
\emph{tiraññeṭuppŭ pracāraṇaṁ innale avasāniccu.}
\end{quote}

\begin{quote}
The election campaign ended yesterday.
\end{quote}

Campaigning is \ml{പ്രചാരണം} \ml{നടത്തുക}.

\begin{cousinweb}
\textbf{\ml{പ്രചാരണം}} (\emph{pracāraṇaṁ}) is a learned word built on Sanskrit \emph{pracāra}, \textquotedblleft{}coming forth, showing oneself; currency, usage\textquotedblright{} (Monier-Williams), with the Sanskrit action-noun ending \emph{-aṇa}: a campaign makes an idea current. Monier-Williams lists \emph{pracāraṇa} itself only as \textquotedblleft{}scattering, strewing\textquotedblright{}, so the campaign sense is a modern one.
\end{cousinweb}

\begin{grammarlens}[title={What you've built}]
One more word for giving reasons and opinions.
\end{grammarlens}

\subsection*{Your turn}
\begin{itemize}
\raggedright
  \item \emph{Say it:} \emph{pracāraṇaṁ}
  \item \emph{Say it:} \emph{pracāraṇaṁ}, once more
  \item \emph{Say it:} \emph{tiraññeṭuppŭ pracāraṇaṁ innale avasāniccu.} — The election campaign ended yesterday.
  \item \emph{Recall:} say the Malayalam for a candidate, then the Malayalam for a ballot, then say \emph{pracāraṇaṁ} again
\end{itemize}

\begin{tcolorbox}[breakable,colback=teal!4,colframe=teal!35!black,title={Before you move on}]
What does \ml{പ്രചാരണം} mean? (\textquotedblleft{}A campaign\textquotedblright{}.) Say it once more.
\end{tcolorbox}
\section[\texorpdfstring{\ml{പ്രകടനപത്രിക} (prakaṭanapatrika)}{പ്രകടനപത്രിക (prakaṭanapatrika)}]{\ml{പ്രകടനപത്രിക} (prakaṭanapatrika) — a manifesto}

\label{lesson:ML-C420-prakatanapatrika}



\begin{quote}
Before the new one: say the Malayalam for a ballot, then the Malayalam for a campaign.
\end{quote}

\subsection*{You'll want to know: \ml{പ്രകടനപത്രിക}}
\textbf{\ml{പ്രകടനപത്രിക}} — \emph{prakaṭanapatrika} — \textquotedblleft{}a manifesto\textquotedblright{}.

\begin{quote}
\ml{ഓരോ} \ml{പാർട്ടിയും} \ml{പ്രകടനപത്രിക} \ml{പുറത്തിറക്കി}.
\end{quote}

\begin{quote}
\emph{ōrō pārṭṭiyuṁ prakaṭanapatrika puṟattiṟakki.}
\end{quote}

\begin{quote}
Each party released a manifesto.
\end{quote}

\ml{പ്രകടനം} alone is a show, a demonstration.

\begin{cousinweb}
\textbf{\ml{പ്രകടനപത്രിക}} (\emph{prakaṭanapatrika}) is a Malayalam compound built from Sanskrit \emph{prakaṭana}, \textquotedblleft{}bringing to light\textquotedblright{} (CDIAL 8429), and \emph{pattrikā}, \textquotedblleft{}a leaf; a letter, a document\textquotedblright{} (Monier-Williams): a paper that brings plans to light.
\end{cousinweb}

\begin{grammarlens}[title={What you've built}]
One more word for giving reasons and opinions.
\end{grammarlens}

\subsection*{Your turn}
\begin{itemize}
\raggedright
  \item \emph{Say it:} \emph{prakaṭanapatrika}
  \item \emph{Say it:} \emph{prakaṭanapatrika}, once more
  \item \emph{Say it:} \emph{ōrō pārṭṭiyuṁ prakaṭanapatrika puṟattiṟakki.} — Each party released a manifesto.
  \item \emph{Recall:} say the Malayalam for a ballot, then the Malayalam for a campaign, then say \emph{prakaṭanapatrika} again
\end{itemize}

\begin{tcolorbox}[breakable,colback=teal!4,colframe=teal!35!black,title={Before you move on}]
What does \ml{പ്രകടനപത്രിക} mean? (\textquotedblleft{}A manifesto\textquotedblright{}.) Say it once more.
\end{tcolorbox}
\section[\texorpdfstring{\ml{ഭൂരിപക്ഷം} (bhūripakṣaṁ)}{ഭൂരിപക്ഷം (bhūripakṣaṁ)}]{\ml{ഭൂരിപക്ഷം} (bhūripakṣaṁ) — a majority}

\label{lesson:ML-C420-bhuripaksam}



\begin{quote}
Before the new one: say the Malayalam for a campaign, then the Malayalam for a manifesto.
\end{quote}

\subsection*{You'll want to know: \ml{ഭൂരിപക്ഷം}}
\textbf{\ml{ഭൂരിപക്ഷം}} — \emph{bhūripakṣaṁ} — \textquotedblleft{}a majority\textquotedblright{}.

\begin{quote}
\ml{അവർക്ക്} \ml{വലിയ} \ml{ഭൂരിപക്ഷം} \ml{കിട്ടി}.
\end{quote}

\begin{quote}
\emph{avarkkŭ valiya bhūripakṣaṁ kiṭṭi.}
\end{quote}

\begin{quote}
They got a big majority.
\end{quote}

An absolute majority is \ml{കേവല} \ml{ഭൂരിപക്ഷം}.

\begin{cousinweb}
\textbf{\ml{ഭൂരിപക്ഷം}} (\emph{bhūripakṣaṁ}) is built from Sanskrit \emph{bhūri}, \textquotedblleft{}much, many\textquotedblright{} (Monier-Williams), and \emph{pakṣa}, \textquotedblleft{}a wing, a side\textquotedblright{} (CDIAL 7627): the side with many.
\end{cousinweb}

\begin{grammarlens}[title={What you've built}]
One more word for giving reasons and opinions.
\end{grammarlens}

\subsection*{Your turn}
\begin{itemize}
\raggedright
  \item \emph{Say it:} \emph{bhūripakṣaṁ}
  \item \emph{Say it:} \emph{bhūripakṣaṁ}, once more
  \item \emph{Say it:} \emph{avarkkŭ valiya bhūripakṣaṁ kiṭṭi.} — They got a big majority.
  \item \emph{Recall:} say the Malayalam for a campaign, then the Malayalam for a manifesto, then say \emph{bhūripakṣaṁ} again
\end{itemize}

\begin{tcolorbox}[breakable,colback=teal!4,colframe=teal!35!black,title={Before you move on}]
What does \ml{ഭൂരിപക്ഷം} mean? (\textquotedblleft{}A majority\textquotedblright{}.) Say it once more.
\end{tcolorbox}
\section[\texorpdfstring{\ml{ന്യൂനപക്ഷം} (nyūnapakṣaṁ)}{ന്യൂനപക്ഷം (nyūnapakṣaṁ)}]{\ml{ന്യൂനപക്ഷം} (nyūnapakṣaṁ) — a minority}

\label{lesson:ML-C420-nyunapaksam}



\begin{quote}
Before the new one: say the Malayalam for a manifesto, then the Malayalam for a majority.
\end{quote}

\subsection*{You'll want to know: \ml{ന്യൂനപക്ഷം}}
\textbf{\ml{ന്യൂനപക്ഷം}} — \emph{nyūnapakṣaṁ} — \textquotedblleft{}a minority\textquotedblright{}.

\begin{quote}
\ml{ആ} \ml{ഗ്രാമത്തിൽ} \ml{അവർ} \ml{ന്യൂനപക്ഷമാണ്}.
\end{quote}

\begin{quote}
\emph{ā grāmattil avar nyūnapakṣamāṇŭ.}
\end{quote}

\begin{quote}
In that village they are a minority.
\end{quote}

Plural with -\ml{ങ്ങൾ}: \ml{ന്യൂനപക്ഷങ്ങൾ}, minorities.

\begin{cousinweb}
\textbf{\ml{ന്യൂനപക്ഷം}} (\emph{nyūnapakṣaṁ}) is built from Sanskrit \emph{nyūna}, \textquotedblleft{}less, deficient\textquotedblright{} (Monier-Williams), and \emph{pakṣa}, \textquotedblleft{}a side\textquotedblright{} (CDIAL 7627): the smaller side.
\end{cousinweb}

\begin{grammarlens}[title={What you've built}]
One more word for giving reasons and opinions.
\end{grammarlens}

\subsection*{Your turn}
\begin{itemize}
\raggedright
  \item \emph{Say it:} \emph{nyūnapakṣaṁ}
  \item \emph{Say it:} \emph{nyūnapakṣaṁ}, once more
  \item \emph{Say it:} \emph{ā grāmattil avar nyūnapakṣamāṇŭ.} — In that village they are a minority.
  \item \emph{Recall:} say the Malayalam for a manifesto, then the Malayalam for a majority, then say \emph{nyūnapakṣaṁ} again
\end{itemize}

\begin{tcolorbox}[breakable,colback=teal!4,colframe=teal!35!black,title={Before you move on}]
What does \ml{ന്യൂനപക്ഷം} mean? (\textquotedblleft{}A minority\textquotedblright{}.) Say it once more.
\end{tcolorbox}
\section[\texorpdfstring{\ml{വിജയിക്കുക} (vijayikkuka)}{വിജയിക്കുക (vijayikkuka)}]{\ml{വിജയിക്കുക} (vijayikkuka) — to win; to succeed}

\label{lesson:ML-C420-vijayikkuka}



\begin{quote}
Before the new one: say the Malayalam for a majority, then the Malayalam for a minority.
\end{quote}

\subsection*{You'll want to know: \ml{വിജയിക്കുക}}
\textbf{\ml{വിജയിക്കുക}} — \emph{vijayikkuka} — \textquotedblleft{}to win; to succeed\textquotedblright{}.

\begin{quote}
\ml{അവൾ} \ml{ആയിരം} \ml{വോട്ടിന്} \ml{വിജയിച്ചു}.
\end{quote}

\begin{quote}
\emph{avaḷ āyiraṁ vōṭṭinŭ vijayiccu.}
\end{quote}

\begin{quote}
She won by a thousand votes.
\end{quote}

Past: \ml{വിജയിച്ചു}, \emph{vijayiccu}.

The margin takes -\ml{ന്}: \ml{ആയിരം} \ml{വോട്ടിന്}. The noun is \ml{വിജയം}, a victory.

\begin{cousinweb}
\textbf{\ml{വിജയിക്കുക}} (\emph{vijayikkuka}) is built from Sanskrit \emph{vijaya}, \textquotedblleft{}victory, conquest\textquotedblright{} (Monier-Williams), with the Malayalam verb ending -\ml{ിക്കുക}.
\end{cousinweb}

\begin{grammarlens}[title={What you've built}]
One more word for giving reasons and opinions.
\end{grammarlens}

\subsection*{Your turn}
\begin{itemize}
\raggedright
  \item \emph{Say it:} \emph{vijayikkuka}
  \item \emph{Say it:} \emph{vijayikkuka}, once more
  \item \emph{Say it:} \emph{avaḷ āyiraṁ vōṭṭinŭ vijayiccu.} — She won by a thousand votes.
  \item \emph{Recall:} say the Malayalam for a majority, then the Malayalam for a minority, then say \emph{vijayikkuka} again
\end{itemize}

\begin{tcolorbox}[breakable,colback=teal!4,colframe=teal!35!black,title={Before you move on}]
What does \ml{വിജയിക്കുക} mean? (\textquotedblleft{}To win; to succeed\textquotedblright{}.) Say it once more.
\end{tcolorbox}
\section[(dialogue)]{Words from chapters 413–420 — a review}

\label{lesson:ML-R420-review-413-420}



\begin{quote}
Say the Malayalam for to recommend and to remark.
\end{quote}

\subsection*{Your turn}
Say these aloud:

\begin{itemize}
\raggedright
  \item to recommend, to remark, to commend, to urge, to express
  \item a principle, an understanding, a theory, a source, per cent
  \item society, the general public, population, equality, freedom
  \item justice, injustice, unfair, human rights, discrimination
  \item a regulation, to abide by, to obey, legal, illegal
  \item punishment, to punish, a ban, to prohibit, to allow
  \item an election, a vote, a voter, a candidate, a ballot
  \item a campaign, a manifesto, a majority, a minority, to win
\end{itemize}

\begin{tcolorbox}[breakable,colback=teal!4,colframe=teal!35!black,title={Before you move on}]
To recommend? (\textbf{\ml{ശുപാർശ} \ml{ചെയ്യുക}}, \emph{śupārśa ceyyuka}.) To remark? (\textbf{\ml{അഭിപ്രായപ്പെടുക}}, \emph{abhiprāyappeṭuka}.) To commend? (\textbf{\ml{പ്രശംസിക്കുക}}, \emph{praśaṁsikkuka}.) To urge? (\textbf{\ml{പ്രേരിപ്പിക്കുക}}, \emph{prērippikkuka}.) To express? (\textbf{\ml{പ്രകടിപ്പിക്കുക}}, \emph{prakaṭippikkuka}.) A principle? (\textbf{\ml{തത്വം}}, \emph{tatvaṁ}.) An understanding? (\textbf{\ml{ധാരണ}}, \emph{dhāraṇa}.) A theory? (\textbf{\ml{സിദ്ധാന്തം}}, \emph{siddhāntaṁ}.) A source? (\textbf{\ml{ഉറവിടം}}, \emph{uṟaviṭaṁ}.) Per cent? (\textbf{\ml{ശതമാനം}}, \emph{śatamānaṁ}.) Society? (\textbf{\ml{സമൂഹം}}, \emph{samūhaṁ}.) The general public? (\textbf{\ml{പൊതുജനം}}, \emph{potujanaṁ}.) Population? (\textbf{\ml{ജനസംഖ്യ}}, \emph{janasaṁkhya}.) Equality? (\textbf{\ml{സമത്വം}}, \emph{samatvaṁ}.) Freedom? (\textbf{\ml{സ്വാതന്ത്ര്യം}}, \emph{svātantryaṁ}.) Justice? (\textbf{\ml{നീതി}}, \emph{nīti}.) Injustice? (\textbf{\ml{അനീതി}}, \emph{anīti}.) Unfair? (\textbf{\ml{അന്യായമായ}}, \emph{anyāyamāya}.) Human rights? (\textbf{\ml{മനുഷ്യാവകാശം}}, \emph{manuṣyāvakāśaṁ}.) Discrimination? (\textbf{\ml{വിവേചനം}}, \emph{vivēcanaṁ}.) A regulation? (\textbf{\ml{ചട്ടം}}, \emph{caṭṭaṁ}.) To abide by? (\textbf{\ml{പാലിക്കുക}}, \emph{pālikkuka}.) To obey? (\textbf{\ml{അനുസരിക്കുക}}, \emph{anusarikkuka}.) Legal? (\textbf{\ml{നിയമപരമായ}}, \emph{niyamaparamāya}.) Illegal? (\textbf{\ml{നിയമവിരുദ്ധമായ}}, \emph{niyamaviruddhamāya}.) Punishment? (\textbf{\ml{ശിക്ഷ}}, \emph{śikṣa}.) To punish? (\textbf{\ml{ശിക്ഷിക്കുക}}, \emph{śikṣikkuka}.) A ban? (\textbf{\ml{വിലക്ക്}}, \emph{vilakkŭ}.) To prohibit? (\textbf{\ml{നിരോധിക്കുക}}, \emph{nirōdhikkuka}.) To allow? (\textbf{\ml{അനുവദിക്കുക}}, \emph{anuvadikkuka}.) An election? (\textbf{\ml{തിരഞ്ഞെടുപ്പ്}}, \emph{tiraññeṭuppŭ}.) A vote? (\textbf{\ml{വോട്ട്}}, \emph{vōṭṭŭ}.) A voter? (\textbf{\ml{വോട്ടർ}}, \emph{vōṭṭar}.) A candidate? (\textbf{\ml{സ്ഥാനാർത്ഥി}}, \emph{sthānārtthi}.) A ballot? (\textbf{\ml{ബാലറ്റ്}}, \emph{bālaṟṟŭ}.) A campaign? (\textbf{\ml{പ്രചാരണം}}, \emph{pracāraṇaṁ}.) A manifesto? (\textbf{\ml{പ്രകടനപത്രിക}}, \emph{prakaṭanapatrika}.) A majority? (\textbf{\ml{ഭൂരിപക്ഷം}}, \emph{bhūripakṣaṁ}.) A minority? (\textbf{\ml{ന്യൂനപക്ഷം}}, \emph{nyūnapakṣaṁ}.) To win? (\textbf{\ml{വിജയിക്കുക}}, \emph{vijayikkuka}.)
\end{tcolorbox}
